Los resultados confirman que la notación asintótica predice correctamente la tendencia general de los algoritmos. Los algoritmos de ordenamiento siguen $O(n\log n)$ (Figura~\ref{fig:sort_complejidad}), mientras que naive y Strassen se ajustan a $O(n^3)$ y $O(n^{2.807})$ respectivamente (Figura~\ref{fig:matrix_complejidad}). Sin embargo, la complejidad temporal por sí sola no explica todo el comportamiento observado. Por ejemplo, en memoria, Quick Sort y
\texttt{std::sort} se mantienen prácticamente constantes por ser \textit{in-place}, frente al crecimiento notoriamente mayor de Patience Sort (Figura~\ref{fig:sort_memoria}); similarmente, Strassen utiliza notablemernte más memoria que naive pese a su ventaja en tiempo (Figuras~\ref{fig:matrix_dominio} y~\ref{fig:matrix_memoria}). Esto evidencia que algoritmos con complejidad temporal similar o mejor pueden diferir considerablemente en su uso de memoria, un trade-off que solo el análisis experimental permite capturar.